\subsection{Обучение и функции потерь}
Модель обучалась с teacher forcing и кросс-энтропией следующего токена.
Для сравнения использовалась token-normalized validation loss:
\[
 L_{\mathrm{token}}=
 -\frac{\sum_{i}\sum_{t\in T_i}\log p_\theta(y_{it}\mid x_i,y_{i,<t})}
 {\sum_i |T_i|},
\]
где $T_i$ содержит целевые токены без padding. Обычный \code{eval_loss}
Trainer усредняется иначе и не подменяет эту величину.
Train loss в графиках --- сохранённое среднее по интервалу логирования;
поэтому его значения не следует трактовать как ту же оценку, что
validation loss, рассчитанный по всем целевым токенам.

В основных поздних опытах использовались AdamW, weight decay 0.01,
clipping 1.0, warmup первых 5\% шагов и cosine decay.
Случайная инициализация и выборки закреплены seed 42.
Смешанная точность выбиралась по GPU: BF16 на части запусков A100/H100,
FP16 на T4 и в итоговом A100-запуске. Конкретные LR, batch и бюджет
указаны у соответствующих экспериментов.

При FP16 GradScaler может пропустить обновление и уменьшить scale.
Поэтому для новых запусков различаются плановые попытки шага Trainer
и фактически успешные обновления весов. Обработанный overflow при конечном
loss и продолжающемся обучении не равнозначен аварийному завершению.

\subsection{Локальный \pass{} и качество кода}
Для каждого условия генерируется ровно один итоговый кандидат,
в том числе при beam search. Эталонный и сгенерированный
\code{solve()} исполняются ограниченным локальным evaluator:
timeout 2 секунды, сравнение чисел с \code{rtol=atol=1e-6}.
Если эталон исполнился и вернул допустимый результат, пример входит
в знаменатель:
\[
 \widehat{\mathrm{pass@1}}=
 \frac{\sum_{i\in E}\mathbf{1}[\widehat{r}_i\simeq r_i]}{|E|}.
\]
Ошибка предсказанной программы при допустимом эталоне считается неуспехом.
Ошибочные эталоны исключаются одинаково у сравниваемых моделей.
Доля ошибок сгенерированного кода считается от всех выбранных строк,
поэтому её знаменатель может отличаться от знаменателя \pass{}.

В Val-20k исполнимы 19\,886 эталонов: исключены 27 неподдерживаемых
программ, 76 runtime errors и 11 недопустимых результатов.
Для Test-001 допустимы 9929, для Test-002 --- 9946 эталонов.
Точное совпадение вопросов, кода, порядка и результатов эталонов
проверялось перед парными сравнениями.

Эта метрика измеряет согласие с выполнением сохранённого эталона,
а не доказанную корректность решения по условию. Ошибка в эталоне
может наказать содержательно правильную программу.

\subsection{Генерация и воспроизводимость}
Использованы KV-cache, батчевая генерация, группировка входов по длине
и восстановление исходного порядка ответов. Основной поздний протокол:
beam 4, length penalty 1.0, до 256 новых токенов, FP16, batch 64.
Обучающий dataset и настройки генерации фиксировались отдельно.
Сохранялись manifests, конфигурации, история loss, checkpoint и
записи исполнения; идентичность файлов проверялась SHA256.

Время обучения, генерации и исполнения кода учитывается раздельно.
Измерение генерации 20\,000 ответов на A100 или T4 не является
проверкой официального лимита на 1919 задачах. Короткие training benchmarks
также не являются оценками скорости авторегрессионного инференса.

